\BeleskeID{04-s033}
% element: 04-s033:e01 — Razlog uvođenja oba komplementa radi prikaza i aritmetike
% element: 04-s033:e02 — Definicija komplementa maksimalne vrednosti za r i n
% element: 04-s033:e03 — Binarni prvi komplement i invertovanje svih bita sa formulom
% element: 04-s033:e04 — Decimalni primer: tri cifre, 1000 kodova, 0–499, −123=876, podela vodećih cifara
% element: 04-s033:d01

U osnovi $r$ najveća cifra je $r-1$. Zapis maksimuma $r^n-1$ sastoji se od te cifre na svih $n$ mesta. Oduzimanjem cifre $c_i$ od $r-1$ na svakom mestu nastaje cifra $(r-1)-c_i$ bez pozajmice. U binarnom sistemu to su zamene $0\leftrightarrow1$.

Znak $\equiv$ u izrazu za negativni broj označava način kodovanja, pa ne tvrdi da su negativna matematička vrednost i pozitivna neoznačena vrednost koda jednake. Kod $876$ u tri decimalne cifre prvog komplementa tumačimo kao $876-999=-123$. Kod $999$ predstavlja drugi zapis nule, dok je $000$ njen nenegativni zapis.
